
\section{Limitations and Future Work} \label{sec:limitations-future-work}

\textbf{Correlations between various code tasks}: While our benchmark serves as an interesting lens to analyze code LMs, one might object that output prediction can simply be done with a Python interpreter and that input prediction abilities can be greatly enhanced by equipping a LM with an interpreter, like in GPT-4 Code Interpreter mode. While this is true, we believe that a good code LM still ought to have good code understanding and execution capabilities, similar to that of a strong programmer. We see that base models have a reasonably strong correlation between HumanEval, input prediction, and output prediction score. An interesting future direction is to more deeply investigate the correlations between performance on various code-related tasks such as code completion, execution, bug-finding, and code summarization.

\textbf{Distilling future execution benchmarks}: Our benchmark only measures the input and output prediction accuracy of relatively simple and self-contained Python functions distilled from a single model (Code Llama 34B). It would also be interesting to measure these capabilities on longer and more difficult code snippets, open-domain code samples, or code in other programming languages. As our distillation technique is relatively general, we welcome others to create their own benchmarks measuring the execution of code snippets from other distributions.

\textbf{Variation due to prompt and temperature}: The accuracy of a model on our benchmark may be very sensitive to the prompt and task format \citep{mizrahi2023state}. We try our best to address this by using prompts that are similar as possible across models (see Appendix \ref{sec:appendix-direct-prompts} and \ref{sec:appendix-cot-prompts}) but understand that some prompts may improve the performance of certain models while decrease the performance on others. There are also countless prompting techniques (see \citep{liu2023pre} for a comprehensive survey) that can be tried to improve the performance. We also run all our experiments with $T=0.2$ and $T=0.8$ due to budget constraints, but different temperatures will lead to different performance for all models. One must always be cautious and critical when using benchmarks to compare models. For example, for input prediction, while Phind v2's 47.9\% pass@1 may seem to beat CodeLlama's 46.5\%, the standard deviations of both models with respect to the $800$ samples selected turns out to be around 1.5\%, so this conclusion cannot be made.

\textbf{Information loss due to pass@1}: While the average pass@k metric is common in the code generation literature, it compresses a large amount of information into one number. While we suggest reporting pass@1 and pass@5 for our benchmark, we comment that pass@k is only one perspective of measuring execution ability. We try to shed more light on behaviour by including a bit more analysis throughout this work, but encourage the development of different evaluation and analysis techniques.

\textbf{Fine-tuning}: In our first fine-tuning experiment, we only check for exact string match when decontaminating the fine-tuning set, so there may still be semantic duplication or similar programs with small modifications, which may lead to a higher performance than if those examples were removed. In this work, we only consider the most direct and straightforward fine-tuning scheme. We believe there is room for improvement via more sophisticated techniques, such as using process supervision \citep{uesato2022solving}, fine-tuning on correct CoT generations, or fine-tuning on snippets of code while including the program state after each step. Seeing that models like Phi, WizardCoder, and Phind outperformed Code Llama on HumanEval but not \benchmark inspires the need for a deeper investigation of the utility of finetuning on distilled data from a more powerful model. Lastly, it remains a curiosity whether fine-tuning on execution information can help code generation abilities.

\textbf{Jointly improving code generation and code execution}: As we discovered, distilled models like Phi, Phind, and WizardCoder that are fine-tuned on code generation do not improve significantly on \benchmark compared to their base models. It is unknown whether the opposite is true: does improved fine-tuning on code execution lead to better code generation abilities? It would also be interesting to explore techniques that can lead to improved performance on both code generation and code execution simultaneously.

\textbf{Understanding reasoning from the lens of code}: As future work, we believe that our benchmark serves as a good starting point towards understanding the code reasoning abilities of LM. Many further execution evaluations may be possible, such as testing execution of recursive functions, execution from a natural language description and an input, or execution of a composition of two functions. We find that output prediction serves as a good testbed for understanding CoT failures, because each step clearly corresponds to an operation with a ground truth, so reasoning failures can be pinpointed. We observed many examples of CoT failures due to simple mistakes that the model seems to have knowledge about (see Appendix \ref{sec:appendix-gpt4-cot-output} for examples), and it should be possible to analyze and characterize this behaviour more systematically. 

\textbf{Self-repair}: Lately, self-repair has been used to improve the reasoning and programming abilities of LLMs \citep{chen2023teaching, olausson2023demystifying, madaan2023self, peng2023check, zhang2023self, tyen2023llms}. From our qualitative analysis, we find that when using CoT, many output prediction failures are recitation errors of information the model may already understand. Therefore, we believe that these mistakes may be easier to repair than when the correct reasoning path is not found in the first place, and that \benchmark may be a simpler task to better understand model repair capabilities.